
\begin{abstract}

A centrally symmetric regular polytope determines a product of dihedral groups indexed by its antipodal vertex pairs. We give a complete presentation of the resulting semidirect product using one relation for each orbit of unordered vertex pairs. An explicit index-two rewriting gives complete presentations for all its endpoint halvings. The bipartite halving argument is proved here, including residual connectivity, flag transitivity and the full type-preserving group. A final observation detects an infinite group whenever a pair orbit is omitted from the presentation. The extension construction, halving operation and Schreier method are classical. The purpose here is to put their exact interfaces in one place, including the converse needed to audit claimed defining presentations.

\end{abstract}

\section{Conventions and retained geometric input}\label{conventions-and-retained-geometric-input}

Let \(P\) be a finite regular lattice polytope of rank \(n\ge2\), with distinguished involutions \(t_0,\ldots,t_{n-1}\), group \(L\), and base vertex \(v\). We assume that its first Schläfli entry is at least three. Write \(L_0=\langle t_1,\ldots,t_{n-1}\rangle\), so the vertex set is the set of right cosets of \(L_0\). Superscripts denote right actions and conjugation: \(x^g=g^{-1}xg\). Let \(\alpha\) be a central involution of \(L\) without fixed vertices; write \(v^*=v^\alpha\). There are \(q=|V(P)|/2\) antipodal pairs.

We retain the matching extension theorem of McMullen--Schulte, \emph{Abstract Regular Polytopes} (2002), Theorem 8C5, and its lattice assertion on p.~264, in the form quoted as Theorem 5 and Lemma 7 in Piedade (2023). Precisely, the group and marking constructed below are the full group and marking of a regular lattice polytope \(E_s(P)\) of type \(\{4,p_1,\ldots,p_{n-1}\}\), for each integer \(s\ge2\). Its vertices are identified with the matching Coxeter group, its vertex figures with \(P\), and its standard face sections with the corresponding parabolic sections. We do not present this retained extension theorem as a new result, or claim to have newly read the original book. The group normal form and presentations below are proved directly.

The geometric proof in \S4 specializes the lattice-polytope halving construction of Montero--Weiss. It reuses the previously established graph proof, but includes every needed geometry and group argument here. No earlier project manuscript is a theorem premise of this paper. The orbital-presentation and parity arguments are classical; the companion application attributes the source families and distinguishes their recovery from the corrected 120-cell presentation.

\section{Presentations indexed by pair orbits}\label{presentations-indexed-by-pair-orbits}

The algebraic statement in this section only requires a transitive action of a presented group \(L\) on a set \(V\), with basepoint stabilizer generated by the specified \(t_j\), \(j\ge1\). Finiteness, the lattice condition and central symmetry are unnecessary for Theorem 1. Thus it also supplies the full group-presentation scope of source Corollary 6 beyond the finite spherical applications.

For each unordered pair \(\{x,y\}\) of distinct vertices choose a symmetric label \(m_{xy}\in\{2,3,\ldots\}\cup\{\infty\}\), invariant under \(L\). Let
\[
 W(m)=\langle u_x\ (x\in V)\mid u_x^2=1,\ (u_xu_y)^{m_{xy}}=1\text{ when }m_{xy}<\infty\rangle.
\]
The permutation action of \(L\) on these generators gives an action by automorphisms. Define \(E(m)=W(m)\rtimes L\), with \(u_x^g=u_{x^g}\), and put \(a=u_v\).

\textbf{Theorem 1 (complete orbital presentation).} A defining presentation of \(E(m)\) consists of a defining presentation of \(L\), the relations \(a^2=1\) and \([a,t_j]=1\) for \(j\ge1\), and, for one representative \(\{v,v^w\}\) of every finite-labelled unordered-pair orbit,
\[
 (a w^{-1}a w)^{m_{v,v^w}}=1.                 \tag{1}
\]
It is permissible to omit all relations from an orbit by giving that orbit label \(\infty\). Relations for the two orientations of an unordered pair are redundant.

\textbf{Proof.} The displayed generators in \(E(m)\) satisfy the proposed relations and generate: every vertex is \(v^g\), so every \(u_x\) is a conjugate of \(a\). Conversely, in the abstract presented group define \(u_{v^g}=a^g\). This is well defined: if \(L_0g=L_0h\), then \(h=kg\) for \(k\in L_0\), and \(a^h=(a^k)^g=a^g\). The relations centralizing the generators of \(L_0\) centralize all of \(L_0\). Each pair is carried by an element of \(L\) to a representative in (1), possibly reversing its endpoints; conjugation, inversion and interchanging the two involutions therefore supply every relation of \(W(m)\). The resulting maps from \(W(m)\) and \(L\) satisfy the semidirect conjugation rule. They induce a homomorphism from \(E(m)\) back to the presented group, inverse on every displayed generator to the first map. This proves the presentation and the claimed converse. In particular no implicit omitted orbit relation is supplied by the argument. $\square$

For explicit orbit calculations, two pairs $\{v,v^g\}$ and $\{v,v^h\}$ are equivalent if and only if
\[
h\in L_0gL_0\ \cup\ L_0g^{-1}L_0.
\tag{O}
\]
If an element carrying one pair to the other fixes $v$, it lies in $L_0$ and gives the first double coset. If it exchanges the endpoints, reversing the first pair replaces $g$ by $g^{-1}$ and gives the second. Conversely, either double-coset relation supplies the required map, with or without reversing the pair. Thus ordered point-stabilizer orbits may merge when the orientation is forgotten.

In the matching case put \(m_{x,x^*}=s\) and all other labels equal to two. Reflections belonging to different antipodal pairs commute, and each pair generates its usual dihedral group of order \(2s\). The presentations in disjoint pairs consequently give
\[
 W_s=D_s^q,\qquad E_s=W_s\rtimes L,\qquad
 |E_s|=(2s)^q|L|.                              \tag{2}
\]
Here \(D_s\) always has order \(2s\); in particular \(D_2=C_2\times C_2\).

The orbit of polytope edges is already enforced by \((at_0)^4=1\). Indeed
\[
 (at_0)^2=a(a^{t_0}),\qquad
 (at_0)^4=(a a^{t_0})^2.
\]
The endpoints of an edge cannot be antipodal under the stated lattice hypotheses. Indeed, suppose a central involution $\alpha$ carried an endpoint $v$ to the other endpoint $w$. Every element of the vertex stabilizer $L_0$ would then fix both $v$ and $w$, by centrality. Their join is the edge: it lies below the edge and cannot equal either of the two distinct vertices. Hence $L_0$ would fix the edge. This contradicts the distinguished generator $t_1\in L_0$, which fixes the base vertex and changes the base edge. Conjugating the base flag proves the assertion for every edge. One may therefore replace the edge-orbit relation by the initial label four and omit it from the additional relations. This yields Piedade's Corollary 6 with a precise sufficiency proof. The index on the representative \(w\) is an orbit of unordered pairs, not a list of distances or a selected Coxeter-element orbit.

\section{The complete parity presentation}\label{the-complete-parity-presentation}

Suppose a group is presented on involutions \(a,b,c_2,\ldots,c_n\). Assume its Coxeter part is the string with initial label four, so \([a,c_j]=1\). Require every further defining relator to be either entirely on the letters \(b,c_j\), or one of the two forms \((a w^{-1}a w)^e\) and \((aw)^{2e}\) treated below, where \(w\) is a word in \(b,c_j\). These relators all have even \(a\)-parity. The character \(\varepsilon\) sending \(a\) to \(-1\) and every other marker to \(1\) is then well defined and surjective. Put
\[
 d=aba,\quad H=\ker\varepsilon,
 \qquad \theta(b)=d,\quad\theta(c_j)=c_j.
\]
The map \(\theta\) on words means letter substitution; it need not define an endomorphism of the subgroup on \(b,c_j\).

\textbf{Theorem 2 (parity rewriting).} The group \(H\) is generated by \(d,b,c_2,\ldots,c_n\). Its complete Coxeter part is the fork formed by duplicating the first node of the old \(b,c_2,\ldots,c_n\) string, with commuting \(d,b\). Each additional relator
\[
 (a w^{-1}a w)^e=1
\]
contributes precisely
\[
 (\theta(w)^{-1}w)^e=1,                         \tag{3}
\]
and each additional relator \((aw)^{2e}=1\) contributes precisely
\[
 (\theta(w)w)^e=1.                             \tag{4}
\]
In these formulas \(w\) is a word in \(b,c_j\). If the old presentation has further relators entirely on those letters, both that relator and its \(\theta\)-substitution must be retained. There is no finiteness hypothesis.

\textbf{Proof.} Use the Schreier transversal \(\{1,a\}\). The relation \(a^2=1\) removes the extra Schreier generator. The generators arising from the other letters are \(b,aba,c_j,ac_ja=c_j\), as stated. Conjugating a word by \(a\) implements \(\theta\). The involution and commuting relations rewrite to the fork relations, while an old relation on \(b,c_j\) gives itself and its conjugate. The initial relation \((ab)^4\) becomes \((db)^2\). The two rewritten versions of (3) are the relation for \(R=\theta(w)^{-1}w\) and its \(a\)-conjugate \(w^{-1}\theta(w)=R^{-1}\), so they are equivalent. In (4), write \(R=\theta(w)w\). Its conjugate is \(w\theta(w)=wRw^{-1}\), again redundant. These are all Schreier relators, proving completeness. $\square$

For \(E_s\), \(\varepsilon\) is total reflection parity on \(W_s\), extended trivially on \(L\). Thus
\[
 H_s=\ker(W_s\longrightarrow C_2)\rtimes L,
 \qquad |H_s|=(2s)^q|L|/2.                    \tag{5}
\]
When \(w=\alpha\) is an involution in \(L\), both \(w\) and \(\theta(w)=awa\) are involutions in the respective old-subgroup copies. Formula (3) becomes \((\theta(\alpha)\alpha)^s=1\). In applications writing \(\alpha=\beta^h\), this is exactly \((\widetilde\beta^h\beta^h)^s=1\), where \(\widetilde\beta=\theta(\beta)\).

\section{Geometry and actual residues}\label{geometry-and-actual-residues}

\textbf{Lemma 3a (bipartite halving).}
Let $Q$ be a regular abstract polytope of rank $r\ge3$ whose face poset is a lattice and whose vertex--edge graph is bipartite. Suppose its distinguished involutions are $a,b,c_2,\ldots,c_{r-1}$, where $a$ changes the vertex and $b$ changes the edge. Replace the first two types by the two vertex colors, retain every proper face of rank at least two, and declare vertices of opposite colors incident precisely when joined by an edge. Incidence with higher faces is containment. This is a thin, residually connected geometry. Its full type-preserving group is the color-preserving subgroup of $\operatorname{Aut}(Q)$, generated by $aba,b,c_2,\ldots,c_{r-1}$. With these generators it has the intersection property and is isomorphic to its maximal-parabolic coset geometry.

\textbf{Proof.}
An edge has two distinct vertices. Because the face poset is a lattice, it is their join: the join is at most the edge and cannot be either endpoint. Consequently two vertices determine at most one edge, and a face contains that edge if and only if it contains both endpoints. Thus the proposed incidence relation is unambiguous.

We first prove the geometry axiom. The higher members of a flag form a chain. If it contains both vertex colors, their unique edge belongs to every higher member; extend that edge and the higher chain to an old chamber. If just one vertex is given, choose an incident edge inside the smallest higher face, or inside $Q$ when there is no higher member, and then extend. If neither color is given, choose an edge in the smallest higher face, or in $Q$. These choices exist because every such face has rank at least two. The endpoints and an extension of the old chain give a new chamber containing the flag.

For residual connectivity, first fix a vertex $v$ of either color. Its residue is isomorphic to the old vertex residue: a neighbor $w$ corresponds to the unique edge $\{v,w\}$, and higher faces are unchanged. The join observation proves preservation of every incidence. Hence residues of flags containing a new vertex are connected whenever their rank is at least two, by the polytope axioms. For a flag of higher faces only, the common vertex--edge graph is that of its smallest face, or of $Q$ for the empty flag. This graph is connected. Indeed an old chamber gallery projects to a walk among vertices and edges; strong flag connectivity inside that face joins every pair of its vertices. Both vertex colors remain in the new residue. Every other residual object belongs to a residual chamber and therefore meets that connected graph. Its incidence graph is consequently connected.

For thinness, a panel missing a higher type fixes two adjacent vertices, hence fixes their edge. Its two extensions are exactly the two extensions of the corresponding old panel. A panel missing one vertex color fixes the other endpoint and a full chain of higher faces. Its two extensions are the two edges through that endpoint in the selected polygonal two-face. Their other endpoints are distinct by the lattice property. Thus every panel has precisely two extensions.

Choose black for the first vertex type. New chambers are in bijection with old chambers whose distinguished vertex is black: an old edge uniquely supplies its white endpoint. The color-preserving subgroup $K$ acts transitively on these chambers. To see this, an automorphism carrying one selected old chamber to another carries a black vertex to a black vertex; connectedness of the old bipartite graph makes it preserve both color classes. The action is faithful, since fixing all new objects fixes every old vertex and higher face, hence also every old edge. It is free on chambers because the old action is free.

We use the following elementary gallery observation. A finite-rank residually connected geometry has a connected chamber graph, and the chambers containing any fixed flag are connected by moves in its missing types. For a proof, induct on rank. Rank one is immediate. In higher rank, connect two chamber objects by an incidence path; extend each consecutive incident pair to a chamber. Consecutive chosen chambers share an object, and the induction hypothesis in its residue joins them. Applying the same induction within a residue proves the flag version. Thinness now implies that a type-preserving automorphism fixing one chamber fixes its unique adjacent chamber at every panel, hence all chambers and all objects. Since $K$ is already chamber-transitive, it is the whole type-preserving automorphism group.

On the base chamber, $b$ changes the white endpoint, $aba$ changes the black endpoint, and each $c_j$ changes its higher face. The gallery observation shows that these panel involutions generate $K$. More generally, the stabilizer of a base subflag consists exactly of the subgroup generated by its missing-type panel involutions. If $K_I$ denotes the subgroup on a set $I$ of panel types, then $K_I$ fixes the complementary base subflag. Intersecting two such pointwise stabilizers gives
\[
K_I\cap K_J=K_{I\cap J}.
\]
Finally send a coset of a base-object stabilizer to that object's image. This gives a type-preserving bijection onto the constructed geometry. Two incident objects extend to a chamber, so chamber transitivity gives a group element in their two cosets. Conversely an element in both cosets sends the corresponding base objects to an incident pair. Incidence therefore agrees exactly with coset intersection. $\square$

\textbf{Theorem 3 (regular halved geometries).} Under the hypotheses of \S1, the maximal-parabolic coset incidence system of $H_s$, with marking $d,b,c_2,\ldots,c_n$, is a regular hypertope. Its full type-preserving group is $H_s$, with order (5), and its diagram is the fork of Theorem 2.

\textbf{Proof.}
Retain the regular lattice polytope $E_s(P)$ from \S1. Its vertices are $W_s$ and its edges have endpoints $w,u_xw$, in a consistent left-action convention. Reflection parity supplies a bipartition, and these generators give a connected graph. The color character on its full group $W_s\rtimes L$ is precisely $\varepsilon$, so its color-preserving subgroup is $H_s$. Apply Lemma 3a. Its panel generators agree with Theorem 2. The products $db=(ab)^2$ have order two because $ab$ has order four; $dc_j$ is conjugate to $bc_j$, and all other pairs are unchanged. These are the asserted exact diagram labels. $\square$


For the convex spherical inputs, no proper face contains antipodal vertices: a supporting hyperplane has positive distance from the center and cannot contain both \(x\) and \(-x\). If \(F\) is a proper face on \(m\) vertices, the corresponding lower face of \(E_s(P)\) therefore has group
\[
 C_2^m\rtimes\operatorname{Aut}(F),            \tag{6}
\]
with its Danzer marking. This is an actual standard parabolic: its conjugates of \(a\) are exactly the vertex reflections of \(F\), independent in \(W_s\), and its old markers give \(\operatorname{Aut}(F)\). The face is consequently the marked polytope \(2^F\), not merely a group of the expected order.

In the halved geometry a residue at either new vertex is \(P\). At a higher face corresponding to \(F\), its lower residue is the halving of \(2^F\), and its upper residue is the section above \(F\) in \(P\); these form a direct product. This follows directly from the lower and upper intervals of a polytope face, and the compression changes only the lower interval. If all proper faces of \(P\) are simplices, the first factor is the universal finite \(D_m\) geometry: its marking is the even-sign subgroup \(\{x\in C_2^m:\sum x_i=0\}\rtimes S_m\), with the usual fork reflections. The conventional low-rank cases are \(D_2=A_1\times A_1\), \(D_3=A_3\). The second factor is spherical. Thus these families are locally spherical. This is the generalized-cube residue argument of Montero--Weiss, reworked here and valid for every \(s\) because (6) is independent of \(s\).

\section{Central symmetry and iteration}\label{central-symmetry-and-iteration}

If \(s\) is even, for each antipodal pair \(\{x,x^*\}\) let \(z_x=(u_xu_{x^*})^{s/2}\). This is the nonidentity central half-turn of that dihedral factor, independent of the orientation of the pair. The product \(z\) of these \(q\) elements is a nonidentity central involution of \(W_s\), and \(L\) fixes it by permuting the pairs. Hence \(z\) is central in \(E_s\). It acts without fixed points on the vertex set \(W_s\) by translation. The retained lattice extension theorem and Theorem 3 may therefore be applied again. This proves the iteration claim for every finite input admitted here, without an assumption about a particular vertex drawing. Classical cubical toroids with the even periods specified in the application provide further eligible inputs.

\section{What an omitted orbit does}\label{what-an-omitted-orbit-does}

The completeness direction of Theorem 1 is useful even when the prescribed labels are incorrect.

\textbf{Theorem 4 (an omitted pair orbit).} Suppose a proposed orbital presentation omits the relation for some unordered pair of distinct vertices. Give every omitted pair label \(\infty\), keeping all stated finite labels. Its group \(E(m)\), and its reflection-parity kernel, are infinite.

\textbf{Proof.} By Theorem 1 its kernel on \(L\) is exactly the Coxeter group \(W(m)\). Here is a direct verification that the product at an omitted pair has infinite order. In the real vector space with basis \(e_x\), define the symmetric form
\[
 B(e_x,e_x)=1,\qquad
 B(e_x,e_y)=-\cos(\pi/m_{xy})\ (m_{xy}<\infty),
 \quad B(e_x,e_y)=-1\ (m_{xy}=\infty).
\]
The linear involutions \(r_x(v)=v-2B(v,e_x)e_x\) satisfy every finite Coxeter relation. For a finite label, the two-dimensional span has nonsingular positive Gram matrix, and the ambient space is its direct sum with the common fixed subspace; on that span the two reflections have product of order \(m_{xy}\). This also covers label two. Thus \(u_x\mapsto r_x\) defines a representation of \(W(m)\), without requiring that the global form be nondegenerate or that the representation be faithful.

For an omitted pair \(x,y\), the invariant span \(\langle e_x,e_y\rangle\) has product matrix
\[
 M=\begin{pmatrix}3&-2\\2&-1\end{pmatrix},\qquad
 M^k=\begin{pmatrix}2k+1&-2k\\2k&1-2k\end{pmatrix}
 \quad(k\ge0).                               \tag{7}
\]
The sign convention can conjugate this matrix but does not change its order. Equation (7), proved by induction, shows infinite order over characteristic zero. The corresponding product \(u_xu_y\) has even reflection parity, so the parity kernel is already infinite. $\square$

The reflection representation is the classical Coxeter construction; we use only the elementary relation check just given. The finite-pair census in the application decides whether the hypothesis of Theorem 4 actually holds. A successful finite realization of a presentation would give only a quotient and would not exclude this defect.

\section{Attribution and verification boundary}\label{attribution-and-verification-boundary}

The extension is McMullen--Schulte's; geometric halving is due to Montero--Weiss, and the complete bipartite proof above reuses the established graph argument. The parity calculation is Reidemeister--Schreier rewriting. Orbital presentations and the reflection construction are classical. The new source-specific work is the complete orbit audit and correction. GAP verifies the finite censuses. Lean verifies only the group identities and (7); all presentation completeness and geometric arguments remain written.

\begin{thebibliography}{9}
\bibitem{P} C. A. Piedade, Infinite Families of Hypertopes from Centrally Symmetric Polytopes, \emph{Electronic Journal of Combinatorics} 30(2) (2023), P2.20, doi:10.37236/10392.
\bibitem{MS} P. McMullen and E. Schulte, \emph{Abstract Regular Polytopes}, Encyclopedia of Mathematics and its Applications 92, Cambridge University Press, 2002. The precise retained results are quoted in \cite{P}; independent access to the book is not claimed.
\bibitem{MW} A. Montero and A. I. Weiss, Locally spherical hypertopes from generalised cubes, \emph{The Art of Discrete and Applied Mathematics} (2020). The halving and simplex-face ideas are attributed to this work; the necessary proofs are written here.
\end{thebibliography}
